\section{实践案例与工程治理}

\subsection{机器人长视野操作}

Chelsea Finn 展示的厨房实例由多个 skill 和 safety guardrail 组成。我们将任务分成：1)感知与状态估计，2) skill 执行（如拿起锅），3) verification + recovery。

\begin{knowledgebox}{Kitchen Robot 的监控点}
\begin{itemize}
    \item State check：每步验证 gripper/contact 是否正常；
    \item Skill success flag：低层 skill 完成后上报 `success`；
    \item Recovery template：若 skill 失败，高层马上触发 safe fallback；
    \item Human override：在 high-risk action (cutting/pouring) 前插入 human-in-the-loop。
\end{itemize}
\end{knowledgebox}

\subsection{评估与监控矩阵}

引入 multi-metric scoreboard 以追踪 accuracy、self-consistency、hallucination rate、token budget 等。每次 emergent jump 都要同时检查多条曲线，避免因为 single metric 误判。

\begin{table}[H]
\centering
\begin{tabular}{p{3.5cm}p{5cm}p{3cm}}
\toprule
指标 & 描述 & 工具 \\
\midrule
Self-consistency & 多路径生成的 agreement & CoT voting harness \\
Hallucination rate & 事实性错误率 & human eval + KG check \\
Token budget & 资源消耗 & monitoring hook \\
Instruction drift & Prompt log 版本偏移 & versioned dashboard \\
\bottomrule
\end{tabular}
\caption{分层系统的评估矩阵。}
\end{table}

\begin{warningbox}{忽略评估链路的后果}
\begin{itemize}
    \item 单看 accuracy 可能掩盖 hallucination；
    \item Prompt drift 会让 high-level plan 失效但低层仍然执行；
    \item 缺少 token budget 监控会造成 compute overrun。
\end{itemize}
\end{warningbox}

\subsection{Incident grade-book 与 lessons}

\begin{table}[H]
\centering
\begin{tabular}{p{4cm}p{6cm}}
\toprule
字段 & 内容 \\
\midrule
Trigger & hallucination/jump/goal failure \\
Root cause & prompt drift / data shift / skill bug \\
Resolution & rollback prompt / retrain options / add guardrail \\
Next action & drill + doc update \\
\bottomrule
\end{tabular}
\caption{incident grade-book 的模板。}
\end{table}

\begin{importantbox}{Grade-book 的操作流程}
\begin{itemize}
    \item 24 小时内完成 trigger → analysis → resolution 的时间线；
    \item 把 prompt/data/attention snapshot 存入 `notes/incident/`；
    \item 把 lessons 更新到 `lessons.md`，在 weekly retro 分享。
\end{itemize}
\end{importantbox}

\subsection{本章小结}

工程治理需要把 skill execution、LLM planning、监控矩阵串成闭环，并在每次 incident 中把 prompt log + evaluation snapshot 归档。

